\chapter{Methods}

\section{Overview}

This chapter describes the modelling and evaluation procedures used in the thesis. Three families of models are trained and evaluated on the same temporally split cohort described previously. Clinical baselines use only emergency department (ED) triage--derived tabular inputs. The image model is a DenseNet121 classifier trained on chest radiographs. The multimodal model applies early fusion by concatenating a pooled visual embedding with an encoded tabular representation and passing the result through a feedforward classifier. All supervised training uses the same binary pneumonia target and the predefined training, validation, and test partitions.

The methodological design is intentionally aligned with the thesis objective of evaluating multimodal learning under realistic temporal constraints. In particular, all model families operate on data derived from the same ED-linked cohort, the same pneumonia label policy, and the same patient-level temporal split. This shared setup ensures that performance differences can be interpreted primarily as differences in model family and input modality, rather than as artifacts of inconsistent cohort construction or evaluation design.

Figure~\ref{fig:pipeline_overview} summarizes the end-to-end pipeline used in this thesis, from data sources and cohort construction to model training and evaluation.

\begin{figure}[htbp]
    \centering
    \resizebox{\textwidth}{!}{%
    \begin{tikzpicture}[
        font=\small,
        >=Latex,
        node distance=1.6cm,
        box/.style={
            draw=black!70,
            rounded corners=3pt,
            align=center,
            inner sep=6pt,
            minimum height=1.2cm
        }
    ]
    
    % --- DATA ---
    \node[box, text width=3.2cm] (data) {
    \textbf{Data sources}\\
    MIMIC-CXR-JPG\\
    MIMIC-IV-ED / MIMIC-IV
    };
    
    % --- COHORT ---
    \node[box, text width=4.5cm, right=of data] (cohort) {
    \textbf{Cohort construction}\\
    ED linkage\\
    $t_0 =$ imaging time\\
    CheXpert labels (\texttt{u\_ignore})\\
    patient-level temporal split
    };
    
    % --- MODELS ---
    \node[box, text width=5.5cm, right=of cohort] (models) {
    \textbf{Models}\\
    Image: DenseNet121 (pretrain $\rightarrow$ finetune)\\
    Clinical: Logistic / XGBoost\\
    Multimodal: early fusion (concat)
    };
    
    % --- OUTPUT ---
    \node[box, text width=3.0cm, right=of models] (out) {
    \textbf{Output}\\
    P(pneumonia)
    };
    
    % --- EVAL ---
    \node[box, text width=6.1cm, below=1.5cm of models] (eval) {
    \textbf{Evaluation}\\
    AUROC, AUPRC, threshold metrics\\
    paired patient-level bootstrap\\
    calibration (ECE, MCE, Brier)\\
    decision curves, hard-negative analysis
    };
    
    % --- ARROWS ---
    \draw[->] (data) -- (cohort);
    \draw[->] (cohort) -- (models);
    \draw[->] (models) -- (out);
    \draw[->] (out) |- (eval);
    
    \end{tikzpicture}%
    }
    \caption{Overview of the pneumonia detection pipeline. Data are used to construct a temporally aligned cohort anchored at imaging time ($t_0$). Clinical, image-only, and multimodal models are trained and evaluated using multiple complementary metrics and uncertainty analysis.}
    \label{fig:pipeline_overview}
\end{figure}

The stages in Figure~\ref{fig:pipeline_overview} mirror the data construction described in the previous chapter and the training and evaluation procedures detailed in the subsections below.

\section{Clinical Models}

The clinical branch is designed to provide interpretable tabular baselines based exclusively on information available at or before the imaging time. These models serve two purposes. First, they quantify how much signal is available from structured triage data alone. Second, they provide a reference point against which the image-only and multimodal systems can be compared.

Two tabular baselines share identical feature sets. Numeric inputs include vital signs, missingness indicators, and the binary view flags \texttt{is\_pa} and \texttt{is\_ap}. Categorical inputs include \texttt{gender}, \texttt{race}, and \texttt{arrival\_transport}. Numeric features are cast to floating point with missing values retained for imputation. Missingness indicators are filled to one when absent, while view flags are filled to zero. Categorical variables are normalized to string representations, with unknown categories mapped to a dedicated token. This preprocessing strategy ensures that missingness is represented explicitly rather than silently absorbed into imputed values, which is especially important in clinical data.

The first baseline is logistic regression. It is implemented as a preprocessing pipeline followed by an L2-regularized classifier with the SAGA solver, a maximum of ten thousand iterations, and balanced class weights. Numeric features are median-imputed and standardized, while categorical features are imputed with a constant token and one-hot encoded with unknown categories ignored at transform time. The pipeline is fit on the training split only and applied unchanged to validation and test data. This model is intentionally simple and serves as a strong linear reference for the triage feature space.

The second baseline is a gradient-boosted tree model implemented using XGBoost with native categorical handling. The classifier uses a binary logistic objective with \texttt{n\_estimators=2000}, \texttt{max\_depth=5}, \texttt{learning\_rate=0.03}, \texttt{subsample=0.8}, \texttt{colsample\_bynode=0.8}, \texttt{reg\_lambda=1}, \texttt{tree\_method="hist"}, \texttt{enable\_categorical=True}, \texttt{n\_jobs=-1}, and \texttt{random\_state=42}. Categorical variables are converted to pandas categorical type before fitting. Training proceeds with early stopping after forty rounds without improvement on validation AUPRC. The positive class weight is set to the ratio of negative to positive examples in the training data using \texttt{scale\_pos\_weight}. Relative to logistic regression, this model is intended to capture nonlinear relationships and higher-order interactions among triage variables while remaining computationally lightweight and interpretable at a system level.

Taken together, the two clinical baselines define a tabular-only performance envelope for the task. If the image model substantially exceeds them, then the radiograph clearly contains additional signal. If multimodal learning improves over both the image and clinical baselines, that improvement can more credibly be attributed to complementary cross-modal information rather than weak unimodal modelling.

\section{Image Model}

The image model is based on torchvision DenseNet121 initialized with ImageNet weights. DenseNet121 is selected because it provides a strong and well-established convolutional backbone for chest radiograph analysis while remaining computationally manageable in the local training environment. The classifier head is replaced with a linear layer producing a single logit for binary pneumonia prediction. Input images are resized to a fixed square resolution (default 224 pixels), converted to tensors, and normalized using ImageNet statistics.

The imaging branch is trained in two stages: multilabel pretraining followed by task-specific fine-tuning. This design reflects the fact that radiographic pneumonia is not an isolated visual phenomenon but one component of a broader pathology space. Pretraining on multiple chest radiograph findings encourages the backbone to learn more general thoracic representations before specialization to the binary pneumonia task.

For multilabel pretraining, the classifier head is replaced with a multi-output linear layer corresponding to the selected CheXpert labels. Training uses masked binary cross-entropy with logits, where only observed binary labels contribute to the loss. In practice, only entries with valid observed \(0/1\) supervision contribute to optimization, while uncertain and missing labels are masked out. This masking strategy is important because CheXpert-style supervision contains uncertain and missing entries, and these should not be treated identically to observed negative labels.

Optimization during pretraining is performed using AdamW with separate learning rates for backbone and classifier parameters, weight decay of 0.0001, optional mixed-precision training, and gradient clipping. Data augmentation at this stage includes horizontal flipping, small rotations, and mild Gaussian blur. Learning rate scheduling follows a ReduceLROnPlateau strategy. By default, model selection in pretraining is based on validation masked loss (\texttt{val\_loss}), although alternative validation criteria such as micro-averaged AUPRC are supported in the training code. The default pretraining budget is 30 epochs with early stopping patience of 5 epochs without improvement.

Fine-tuning for pneumonia initializes from ImageNet or an optional multilabel checkpoint, depending on the run configuration. When a pretrained checkpoint is used, backbone weights are loaded while the final classifier head is reinitialized for binary output. The fine-tuning loss is binary cross-entropy with logits, using a positive class weight computed from the training split as the ratio of negative to positive examples. Optimization, checkpointing, and learning rate scheduling follow the same general strategy as in pretraining, but model selection is based on validation AUPRC rather than validation loss. The default fine-tuning budget is 40 epochs with early stopping patience of 10 epochs. Compared with pretraining, data augmentation is lighter; the fine-tuning pipeline uses resize, normalization, and mild Gaussian blur rather than the full pretraining augmentation set.

AUPRC is emphasized because pneumonia prevalence in the labeled benchmark is not extremely low but still sufficiently imbalanced that area under the precision--recall curve provides a useful complement to AUROC for model selection and final evaluation.

This two-stage image-training strategy is intended to produce a strong unimodal radiograph baseline. In the context of this thesis, it also serves as the reference against which multimodal fusion must be judged. A central question is not simply whether multimodal learning works in absolute terms, but whether it improves on an already competitive image representation learned from the same cohort and evaluated under the same temporal constraints.

\section{Multimodal Model}

The multimodal architecture combines a DenseNet121 image encoder with a tabular multilayer perceptron. The design goal is to fuse radiographic and triage information while keeping the architecture simple, interpretable at a high level, and directly comparable to the unimodal baselines.

The image branch applies convolutional feature extraction followed by global average pooling to produce a 1024-dimensional embedding. This dimension is inherited directly from the DenseNet121 backbone. The tabular branch consists of a two-layer feedforward network with batch normalization, ReLU activations, and dropout, producing a 128-dimensional representation. The difference in dimensionality reflects the fact that the raw information content and representational complexity of the image branch are substantially larger than those of the structured triage branch.

Figure~\ref{fig:early_fusion_architecture} illustrates the early-fusion design used in this thesis.

\begin{figure}[htbp]
    \centering
    \resizebox{\textwidth}{!}{%
    \begin{tikzpicture}[
        font=\small,
        >=Latex,
        node distance=1.4cm,
        box/.style={
            draw=black!70,
            rounded corners=3pt,
            align=center,
            minimum height=1.1cm,
            inner sep=6pt
        },
        arrow/.style={->, thick}
    ]

    % Inputs
    \node[box, text width=3.0cm] (imgin) {
        \textbf{Image input}\\
        CXR image\\
        224 $\times$ 224
    };

    \node[box, text width=3.0cm, below=2.0cm of imgin] (tabin) {
        \textbf{Tabular input}\\
        ED triage features\\
        preprocessed matrix
    };

    % Branches
    \node[box, text width=3.7cm, right=2.0cm of imgin] (imgenc) {
        \textbf{Image branch}\\
        DenseNet121 feature extractor\\
        global average pooling
    };

    \node[box, text width=3.7cm, right=2.0cm of tabin] (tabenc) {
        \textbf{Tabular branch}\\
        2-layer MLP\\
        BatchNorm + ReLU + Dropout
    };

    % Embeddings
    \node[box, text width=2.8cm, right=1.7cm of imgenc] (imgemb) {
        \textbf{Image embedding}\\
        1024-d
    };

    \node[box, text width=2.8cm, right=1.7cm of tabenc] (tabemb) {
        \textbf{Tabular embedding}\\
        128-d
    };

    % Fusion
    \node[box, text width=3.4cm, right=2.0cm of $(imgemb)!0.5!(tabemb)$] (concat) {
        \textbf{Concatenation}\\
        1152-d fused vector
    };

    \node[box, text width=4.0cm, right=1.9cm of concat] (fusion) {
        \textbf{Fusion head}\\
        hidden layer 256\\
        BatchNorm + ReLU + Dropout 0.2
    };

    \node[box, text width=2.6cm, right=1.8cm of fusion] (out2) {
        \textbf{Output}\\
        pneumonia probability
    };

    % Arrows
    \draw[arrow] (imgin) -- (imgenc);
    \draw[arrow] (tabin) -- (tabenc);

    \draw[arrow] (imgenc) -- (imgemb);
    \draw[arrow] (tabenc) -- (tabemb);

    \draw[arrow] (imgemb.east) -- ++(0.6,0) |- (concat.west);
    \draw[arrow] (tabemb.east) -- ++(0.6,0) |- (concat.west);

    \draw[arrow] (concat) -- (fusion);
    \draw[arrow] (fusion) -- (out2);

    \end{tikzpicture}%
    }
    \caption{Early-fusion multimodal architecture used in this thesis. A DenseNet121 image branch produces a 1024-dimensional visual embedding, while a two-layer multilayer perceptron encodes tabular triage features into a 128-dimensional representation. The two embeddings are concatenated into a 1152-dimensional fused vector and passed through a fusion head to generate the final pneumonia probability.}
    \label{fig:early_fusion_architecture}
\end{figure}

Early fusion is implemented by concatenating the image and tabular embeddings and passing the result through a fusion head comprising a fully connected layer, normalization, activation, dropout, and a final linear output layer. In the implemented model, the concatenated vector has dimensionality 1152 and is projected through a fusion hidden layer of size 256 before the final binary output. The default dropout rate in both the tabular and fusion branches is 0.2. This architecture is intentionally modest in complexity. Its purpose is not to exhaust the full space of multimodal architectures, but rather to provide a clean and well-defined test of whether triage information adds useful signal when fused with a strong image encoder under realistic timing constraints.

Tabular inputs are preprocessed using a column-wise transformer fitted on the training split, with median imputation and scaling for numeric features and constant imputation plus one-hot encoding for categorical features. This preprocessing is performed independently of the image branch but within the same split discipline as the clinical baselines, thereby preserving comparability and avoiding train--test leakage. The fitted tabular preprocessor is saved as an artifact (\texttt{tabular\_preprocessor.joblib}) to ensure reproducible transformation at validation and test time.

The image backbone may be initialized from ImageNet or a pretrained checkpoint. An optional configuration allows freezing backbone parameters so that only the tabular branch and fusion head are trained. When the backbone is trainable, optimization uses separate parameter groups for backbone and non-backbone components. In the main reported experiments, the multimodal model is trained under the same general evaluation framework as the image-only model so that the performance difference can be interpreted primarily as the effect of adding the tabular modality. The default multimodal training budget is 30 epochs with early stopping patience of 8 epochs, and model selection is based on validation AUPRC. As in the image fine-tuning stage, image preprocessing for multimodal training uses resize, normalization, and mild blur rather than the full augmentation set used during multilabel pretraining.

\section{Training Procedure}

All models are trained on the training partition and evaluated on validation and test partitions without cross-contamination. Neural models use mini-batch training with shuffled training data and deterministic validation and test evaluation. This common split discipline is essential because the thesis focuses not only on absolute performance but also on fair comparison across modalities.

Class imbalance is addressed differently for the three model families. Logistic regression uses balanced class weights. XGBoost uses \texttt{scale\_pos\_weight} computed from the training split. Neural models use weighted binary cross-entropy, with a positive class weight derived from the prevalence of the training labels. These strategies are intended to reduce bias toward the negative class without changing the underlying target definition.

Validation performance guides model selection throughout the system, but the exact criterion differs by training stage. Multilabel pretraining defaults to validation masked loss, whereas image fine-tuning and multimodal training both use validation AUPRC as the primary model-selection criterion. Neural models adjust learning rates when performance plateaus and apply early stopping after a fixed patience. The best validation checkpoint is retained for final test evaluation. Clinical models use validation primarily for monitoring, while XGBoost additionally uses it for early stopping. This consistent reliance on the validation split prevents test-set feedback during model development.

At the implementation level, model outputs, validation histories, checkpoints, and test predictions are saved as artifacts so that later analyses can be run from frozen predictions rather than from repeated retraining. This is particularly important for calibration assessment, decision curve analysis, bootstrap evaluation, and qualitative inspection, all of which are performed on saved outputs to ensure consistency across model families. Saved run artifacts include configuration files, training histories, summary files, serialized models or preprocessors, and prediction CSVs for validation and test partitions.

\section{Evaluation}

Model performance is assessed primarily using AUROC and AUPRC computed from predicted probabilities. Accuracy, precision, recall, and F1 score are additionally reported using a fixed decision threshold of 0.5. All metrics are evaluated on held-out test data with both classes present.

Uncertainty is quantified using patient-level bootstrap resampling. Patients are sampled with replacement, and all associated studies are included in each replicate. Metrics are recomputed for each bootstrap sample, and summary statistics including mean, standard deviation, and percentile-based intervals are reported. This approach respects the grouped structure of the dataset and avoids treating correlated studies from the same patient as independent observations. Although the generic evaluation utilities support a default of 1000 bootstrap replicates, the canonical analyses reported in this thesis use 2000 replicates with a fixed random seed.

Model comparisons use paired bootstrap resampling, ensuring that each replicate evaluates both models on the same resampled patient set. In the implemented evaluation pipeline, prediction files are aligned by \texttt{subject\_id} and, when available, \texttt{study\_id}, and consistency checks are performed to ensure that merged rows correspond to identical targets before paired differences are computed. Differences in AUROC and AUPRC are summarized across replicates. In addition to confidence intervals, the evaluation reports the fraction of bootstrap replicates in which one model outperforms another; this quantity should be interpreted as a bootstrap replicate proportion rather than as a formal frequentist \(p\)-value.

Calibration is evaluated using Brier score and reliability-style binning of predicted probabilities. Expected calibration error (ECE) and maximum calibration error (MCE) are computed using 10 equal-width bins over the unit interval. Bootstrap resampling is also used to summarize uncertainty in calibration metrics. These analyses complement discrimination metrics by assessing whether predicted probabilities align with observed event frequencies.

In addition, decision curve analysis (DCA) is used to assess threshold-dependent net benefit across models. In the canonical evaluation, thresholds range from 0.01 to 0.99 using 99 evenly spaced points. Before DCA curves are computed, model prediction files are checked for row consistency so that threshold-wise comparisons are performed on aligned prediction sets. Decision curve analysis complements AUROC and AUPRC by evaluating whether a model would improve threshold-based decision making relative to simpler strategies such as treating all or treating none.

Finally, qualitative analyses are conducted using gradient-based visualization methods applied to representative cases. These are not treated as substitutes for quantitative evaluation, but as complementary diagnostics that help interpret model behavior in correct and incorrect predictions. All evaluation procedures operate on saved prediction files or saved model checkpoints to ensure consistency, reproducibility, and comparability across experiments.